\documentclass[12pt]{article}
\usepackage[T1]{fontenc}
\usepackage[utf8]{inputenc}
\usepackage{lmodern}
\usepackage{textcomp}
\usepackage{enumitem}
\usepackage{geometry}
\geometry{margin=1in}
\begin{document}
\begin{center}
Answer Key - Constitutional Law - Graduate Level Assessment 24 \
\small (Answers are ordered exactly as in the question paper.)
\end{center}

\section*{Multiple Choice}
\begin{enumerate}[label=\arabic*.]
\item \textbf{Answer:} A
\\ \textit{Options:} A) The law is the only guide for ethical behaviour.; B) The law dictates what moral values should affect our ethical behaviour.; C) Morality plays no role in the concept of law.; D) Moral arguments operate only in hard cases.; E) Morality only applies to specific, complex cases.
\\ \textit{Note:} Correct option: A - The law is the only guide for ethical behaviour.
\item \textbf{Answer:} A
\\ \textit{Options:} A) Dudgeon v UK (1981); B) Marckx v Belgium (1979); C) Osman v UK ( 1998); D) Tyrer v UK (1978); E) McCann v UK (1995)
\\ \textit{Note:} Correct option: A - Dudgeon v UK (1981)
\item \textbf{Answer:} A
\\ \textit{Options:} A) Treaties do no create obligations or rights for third States without their consent; B) Treaties create both obligations and rights for third States; C) Treaties do not create any rights for third States, even when the latter consent.; D) Treaties may create only obligations for third States; E) Treaties create obligations for third States without their consent
\\ \textit{Note:} Correct option: A - Treaties do no create obligations or rights for third States without their consent
\item \textbf{Answer:} C
\\ \textit{Options:} A) Related Doctrine of Merger; B) Contingent remainder; C) Doctrine of Worthier Title; D) Doctrine of Escheat; E) Vested remainder
\\ \textit{Note:} Correct option: C - Doctrine of Worthier Title
\item \textbf{Answer:} B
\\ \textit{Options:} A) His concept of status is misrepresented.; B) It is misinterpreted as a prediction.; C) His idea is considered inapplicable to Western legal systems.; D) It is incorrectly related to economic progression.; E) It is taken literally.
\\ \textit{Note:} Correct option: B - It is misinterpreted as a prediction.
\end{enumerate}

\section*{True / False}
\begin{enumerate}[label=\arabic*.]
\setcounter{enumi}{5}
\item \textbf{Answer:} False
\\ \textit{Options:} A) True; B) False
\\ \textit{Note:} LLM-generated false statement. Correct answer: 'Warranty of Fitness for Particular Purpose.'.
\item \textbf{Answer:} False
\\ \textit{Options:} A) True; B) False
\\ \textit{Note:} LLM-generated false statement. Correct answer: 'The POP will choose wealth over a compassionate society.'.
\item \textbf{Answer:} True
\\ \textit{Options:} A) True; B) False
\\ \textit{Note:} LLM-generated true statement based on MCQ.
\item \textbf{Answer:} True
\\ \textit{Options:} A) True; B) False
\\ \textit{Note:} LLM-generated true statement based on MCQ.
\item \textbf{Answer:} False
\\ \textit{Options:} A) True; B) False
\\ \textit{Note:} LLM-generated false statement. Correct answer: 'It is an agreement between two or more to commit a crime.'.
\end{enumerate}

\section*{Long Form Response}
\begin{enumerate}[label=\arabic*.]
\setcounter{enumi}{10}
\item \textbf{Answer:} Cognitivism.. Explain why this is the correct answer and provide supporting evidence. Reference key facts or concepts that support this choice.
\\ \textit{Note:} Long-form derived from MCQ; award credit for stating the correct idea and giving supporting reasoning.
\item \textbf{Answer:} A bidder may retract his bid at any time until the falling of the hammer.. Explain why this is the correct answer and provide supporting evidence. Reference key facts or concepts that support this choice.
\\ \textit{Note:} Long-form derived from MCQ; award credit for stating the correct idea and giving supporting reasoning.
\end{enumerate}

\vfill
\noindent\textit{End of Answer Key}
\end{document}